\documentclass[11pt,a4paper]{article}
\usepackage[margin=25mm]{geometry}
\usepackage{amsmath,amssymb}
\setlength{\parindent}{0pt}
\setlength{\parskip}{0.6em}
\setlength{\emergencystretch}{2em}
\begin{document}
\begin{center}
  {\LARGE\bfseries M --- My Problem Was Not Written By AI}\\[0.5em]
  {\small Graphs, functional graphs, SCC}
\end{center}

\subsection*{1. Turn the words into edges}

Compress the words into vertices. For each position, if the original
word is $x$ and Frank's word is $y$, we need an edge $x\to y$.

Every occurrence of the same word changes together. So if one source
word needs two different successors, the transformation is impossible.
We output \texttt{You win this time, Frank.}

Ignore self-loops and count repeated edges only once. Every remaining
edge means a word needs to change, so start with the number of these
edges as our answer.

\subsection*{2. The only problem is cycles}

For a chain $a\to b\to c$, we can do $b\to c$ first, then $a\to b$.
So we handle destinations before the words pointing into them.
Cycles are where this gets stuck. Since each vertex has at most one
outgoing edge, SCCs with more than one vertex are exactly the cycles.

A cycle with no outside incoming edge needs a temporary word.
For example, to swap $a$ and $b$, use a fresh word $t$:
\[
  a\to t,\qquad b\to a,\qquad t\to b.
\]
That is three operations for two edges. In general, this costs one
extra operation per cycle.

If there is an outside edge $x\to a$, we can use $x$ instead:
\[
  b\to x,\qquad a\to b,\qquad x\to a.
\]
Both $b$ and $x$ need to become $a$, so merging them is fine.
Now we handle three edges in three operations, with no extra cost.
The same idea works for longer cycles.

\subsection*{3. What to count}

For each SCC with more than one vertex, check whether any edge enters
it from outside. Every cycle vertex already has one incoming edge
from the cycle itself. So there is no outside edge exactly when
every vertex in that SCC has indegree $1$.

So the answer is the number of distinct non-self edges, plus one for
each cycle with no outside incoming edge.

\subsection*{4. Complexity}

After compressing the words, the algorithm takes $O(n)$ time and
$O(n)$ space.

\end{document}
